\section{Preliminaries}

%\lnote{$f^{\oplus}$ should be defined somewhere.}
We use $[n]$ to denote the set $\set{1,2,\ldots, n}$ and \textbf{bold font} (e.g $\bx \sim \mcD$) to denote random variables. For any distribution $\mu$, we use $\mu^k$ to denote $k$-fold the product distribution $\mu \times \cdots \times \mu$.

For any function $f:\bn \to \bits$, we use $f^{\otimes k}: (\bn)^k \to \bits^k$ to denote its $k$-fold direct product,
\begin{equation*}
    f^{\otimes k}(X^{(1)}, \ldots, X^{(k)}) \coloneqq \paren[\big]{f(X^{(1)}), \ldots, f(X^{(k)})}.
\end{equation*}
Similarly, we use $f^{\oplus k}: (\bn)^k \to \bits$ to denote its $k$-fold direct sum,
\begin{equation*}
    f^{\oplus k}(X^{(1)}, \ldots, X^{(k)}) \coloneqq \prod_{i \in [k]}f(X^{(i)}).
\end{equation*}


\begin{definition}[Bernoulli distribution]
    For any $\delta \in [0,1]$, we write $\Ber(\delta)$ to denote the distribution of $\bz$ where $\bz = 1$ with probability $\delta$ and $0$ otherwise.
\end{definition}

\begin{definition}[Binomial distribution]
    For any $k \in \N$, $\delta \in [0,1]$, we write $\Bin(k, \delta)$ to denote the sum of $k$ independent random variables drawn from $\Ber(\delta)$.
\end{definition}

\begin{fact}[Chernoff bound]
    \label{fact:chernoff}
    Let $\bz_1, \ldots, \bz_k$ be independent and each bounded within $[0,1]$ and $\bZ \coloneqq \sum_{i \in [k]}\bz_i$. For any threshold $t \leq \mu \coloneqq \Ex[\bZ]$,
    \begin{equation*}
        \Pr[\bZ \leq t] \leq \exp\paren*{-\tfrac{(\mu - t)^2}{2\mu}}.
    \end{equation*}
    Similar bounds for the probability $\bZ$ exceeds its mean hold. For example,
    \begin{equation*}
        \Pr[\bZ \geq 2\mu] \leq \exp\paren*{-\tfrac{\mu}{3}}.
    \end{equation*}
    Furthermore, the above bounds also hold for any random variable $\bY$ satisfying $\Ex[e^{\lambda \bY}] \leq \Ex[e^{\lambda \bZ}]$ for all $\lambda \in \R$.
\end{fact}




\subsection{Randomized vs.~deterministic decision trees}
We will prove all of our results with respect to the expected depth of a \emph{randomized} decision tree. In this subsection, we formally define \emph{deterministic} and \emph{randomized} decision trees and prove that our results easily extend to the deterministic setting.

\begin{definition}[Deterministic decision tree]
    A deterministic decision tree, $T: \bits^n \to \bits$, is a binary tree with two types of nodes: Internal nodes each query some $x_i$ for $i \in [n]$ and have two children whereas leaf nodes are labeled by a bit $b \in \bits$ and have no children. On input $x \in \bits^n$, $T(x)$ is computed as follows: We proceed through $T$ starting at the root. Whenever at an internal node that queries the $i^{\text{th}}$ coordinate, we proceed to the left child if $x_i = -1$ and right child if $x_i = +1$. Once we reach a leaf, we output the label of that leaf.
\end{definition}
\begin{definition}[Randomized decision tree]
    A randomized decision tree, $\mcT:\bits^n \to \bits$, is distribution over deterministic decision trees. On input $x \in \bits^n$, it first draws $\bT \sim \mcT$ and then outputs $\bT(x)$.
\end{definition}
\begin{definition}[Expected depth]
    For any deterministic decision tree $T: \bits^n \to \bits$ and distribution $\mu$ on $\bits^n$, we use $\overline{\mathrm{Depth}}^{\mu}(T)$ to denote the expected depth of $T$, which is the expected number of coordinates $T$ queries on a random input $\bx \sim \mu$. Similarly, for a randomized decision tree $\mcT$, $\overline{\mathrm{Depth}}^{\mu}(\mcT) \coloneqq \Ex_{\bT \sim \mcT}[\overline{\mathrm{Depth}}^{\mu}(\bT)]$.
\end{definition}



We write $\overline{\mathrm{Depth}}^{\mu}(\cdot,\cdot)$ to denote the minimum expected depth of \emph{any} decision tree, including randomized decision trees. Thus, all of our main results, as written, hold for \emph{randomized} decision trees; however, equivalent statements are true if we restrict ourselves to only \emph{deterministic} decision trees, with only a small change in constants, as easily seen from the following claim.
\begin{claim}
    \label{claim:connect-det-rand}
    For any $f:\bn \to \bits$, distribution $\mu$, and constant $\eps$,
    \begin{equation*}
         \odepthmu(f,2\eps) \leq \odepthmu_{\det}(f, 2\eps) \leq 2 \cdot \odepthmu(f, \eps).
    \end{equation*}
\end{claim}
\begin{proof}
    The left-most inequality follows immediately from that fact that any deterministic decision tree is also a randomized decision tree. For the second inequality, given any randomized decision tree $\mcT$ with error $\eps$ and expected depth $d$, we'll construct a deterministic decision tree $T$ with error at most $2\eps$ and expected depth at most $2d$. First, we decompose the expected error of $\mcT$:
    \begin{equation*}
        \eps = \Prx_{\bx \sim \mu}[\mcT(\bx) \neq f(\bx)] = \Ex_{\bT \sim \mcT}\bracket[\Big]{\Prx_{\bx \sim \mu}[\bT(\bx) \neq f(\bx)]}.
    \end{equation*}
    Applying Markov's inequality, if we sample $\bT \sim \mcT$, with probability at least $1/2$, it has error at most $\eps$. Similarly, since the expected depth of $\mcT$ is $d$, with probability at least $1/2$, $\bT$ will have expected depth at most $2d$. By union bound, there is a nonzero probability that we choose a single (deterministic) tree with error at most $2\eps$ and expected depth at most $2d$.
\end{proof}
\Cref{claim:connect-det-rand} immediately allows direct sum theorems for \emph{randomized} decision trees to also apply to \emph{deterministic} decision trees.
\begin{corollary}[Randomized direct sum theorems imply deterministic ones]
    \label{cor:rand-to-det-DP}
    Suppose that a randomized direct sum theorem of the following form holds. For a function $f: \bn \to \bits$, distribution $\mu$  over $\bn$,  $k\in \mathbb{N}$, and constants $\varepsilon, \delta, M$,
    \[ \overline{\mathrm{Depth}}^{\mu^{\otimes k}}(f^{\otimes k},\varepsilon) \ge M \cdot \overline{\mathrm{Depth}}^{\mu}(f,\delta).\]
    Then, 
     \[ \overline{\mathrm{Depth}}_{\det}^{\mu^{\otimes k}}(f^{\otimes k},\varepsilon) \geq \overline{\mathrm{Depth}}^{\mu^{\otimes k}}(f^{\otimes k},\varepsilon) \ge M \cdot \overline{\mathrm{Depth}}^{\mu}(f,\delta) \geq  \tfrac{M}2 \cdot \overline{\mathrm{Depth}}^{\mu}_{\mathrm{det}}(f,2\delta).\]
\end{corollary}
With \Cref{cor:rand-to-det-DP} in mind, the remainder of this paper will only consider \emph{randomized} decision trees.

% \Cref{thm:main-formal}, as written, holds for \textit{randomized} decision trees; however, an equivalent statement is true if we restrict ourselves to only \textit{deterministic} decision trees. In fact, the strongest version of \Cref{thm:main-formal}, which restricts to deterministic decision trees \textit{only} on the right hand side, also holds. We prove that this stronger version follows immediately from \Cref{thm:main-formal}. Let $\overline{\mathrm{Depth}}_{\mathrm{det}}^{\mu}(f,\varepsilon)$ denote the minimum expected depth of any \textit{deterministic} decision tree that computes $f$ to error $\varepsilon$ w.r.t.~$\mu$. We show:

% \begin{claim}[Randomized direct product theorems imply deterministic ones]
%     Suppose that a randomized direct product theorem holds: for a function $f: \bn \to \bits$, distribution $\mu$  over $\{0,1\}^n$,  $k\in \mathbb{N}$, and constants $\varepsilon, \delta, M$,
%     \[ \overline{\mathrm{Depth}}^{\mu^{\otimes k}}(f^{\otimes k},\varepsilon) \ge M \cdot \overline{\mathrm{Depth}}^{\mu}(f,\delta).\]
%     Then, 
%      \[ \overline{\mathrm{Depth}}^{\mu^{\otimes k}}(f^{\otimes k},\varepsilon) \ge \frac{M}2 \cdot \overline{\mathrm{Depth}}^{\mu}_{\mathrm{det}}(f,2\delta).\]
% \end{claim}

% \begin{proof} It is sufficient to prove that $\overline{\mathrm{Depth}}^{\mu}_{\mathrm{det}}(f,2\delta)\leq 2 \cdot \overline{\mathrm{Depth}}^{\mu}(f,\delta)$. Let $d=\overline{\mathrm{Depth}}^{\mu}(f,\delta)$. Then, there exists a distribution $\mathcal{T}$ over deterministic trees such that 
%     \begin{enumerate}
%         \item $\Ex_{\bT\sim \mathcal{T}}[\overline{\mathrm{Depth}}^{\mu}(\bT)]\leq d$ and
%         \item $\Ex_{\bT\sim \mathcal{T}}[\Pr_{\bx\sim \mu }[\bT(\bx)\neq f(\bx)]]\leq \delta.$
%     \end{enumerate}
%     Applying Markov bounds to both of the above inequalities yields
%     \begin{enumerate}
%         \item $\Pr_{\bT\sim \mathcal{T}}[\overline{\mathrm{Depth}}^{\mu}(\bT)>2d]<1/2$ and
%         \item $\Pr_{\bT\sim \mathcal{T}}[\Pr_{\bx\sim \mu }[\bT(\bx)\neq f(\bx)]>2\delta]<1/2$.
%     \end{enumerate}
%     By a union bound, we conclude that there exists a deterministic decision tree $T$ in the support of $\mathcal{T}$ such that 
%     \begin{enumerate}
%         \item $\overline{\mathrm{Depth}}^{\mu}(T)\leq 2d$ and
%         \item $\Pr_{\bx\sim \mu }[T(\bx)\neq f(\bx)]\leq 2\delta$.
%     \end{enumerate}
%     We conclude that $\overline{\mathrm{Depth}}^{\mu}_{\mathrm{det}}(f,2\delta)\leq 2d$.
% \end{proof}

% Trivially, $ \overline{\mathrm{Depth}}^{\mu^{\otimes k}}_{\mathrm{det}}(f^{\otimes k},\varepsilon)\geq  \overline{\mathrm{Depth}}^{\mu^{\otimes k}}(f^{\otimes k},\varepsilon)$, so the above claim does indeed yield a fully deterministic direct product theorem. 
